% ===========================================================
% 《线性代数9讲》第 8 讲  相似理论 —— 片段 B（本讲中段）
% 源：线代9讲强化.pdf  PDF p95–p98（印刷页 84–87），共 4 页
%
% 处理说明：
%   1) OPD 方法论标记剔除（依 AGENTS.md §7.1 决定 2 的统一口径）：
%      · \fsec 标题里的 OPD 括号连同其中的 OPD 词组**整段删除**：
%        「四、$A$ 与 $B$ 相似」（原稿后附 $O_4$（盯住目标 4）+$D_1$（常规操作）
%          +$D_{21}$（观察研究对象）+$P_4$（逆否思路））——已删；
%      · 正文蓝色纯 OPD 代码 / 方法论语句以 `% OPD 蓝色标注（原稿）：…` 留痕，不排入正文：
%        PDF p97（86）（1）末句右侧 →$D_{22}$（转换等价表述），当一个结论有"充要""充分"
%          "必要""否定"等条件，从而构成了一个好的命题点"时，便给考生指出了重点复习的方向；
%      · 版心左侧竖排模块名「隐含条件体系块」（PDF p97，OPD 体系块标记）→ 删除并留痕；
%      · 原稿右上角蓝色装饰性二维码 1 个（PDF p97，非正文）→ 未转写。
%   2) **数学操作旁注保留**（照原文排入，均为数学内容，非 OPD 代码）：
%      · PDF p95（84）$Q$ 下方蓝色手写箭头标注「$A=Q\Lambda Q^{-1}$」及其下「$Q^{\mathrm T}$」
%        （指 $Q^{-1}=Q^{\mathrm T}$）；
%      · PDF p96（85）三行蓝色手写模板表（$A\,|\,f(A)$、$\lambda\,|\,f(\lambda)$、$\xi\,|\,\xi$）
%        及其下 5 行特征值对应关系；
%      · PDF p96（85）【注】框下方蓝色手写「若 $\lambda_1=100$…你是否看出 $A$ 的主要成分是
%        $\lambda_1\xi_1\xi_1^{\mathrm T}$？」；
%      · PDF p97（86）（2）（3）两问右侧蓝色手写证明（含 $P^{-1}$、$P$ 的上下标注）。
%   3) 页眉（「第8讲 相似理论」）、页脚（84/85/86/87 及云纹）为每页重复装饰，按规范忽略。
%   4) 本片 4 页**无「习题」栏**（已逐页核实 PDF p95–p98）。
%
% 接缝说明：
%   · 本片首行承 PDF p94（印刷页 83）例 8.5 的【解】法一（$P^{-1}$ 的初等变换求法），
%     该例题解答跨页，本片以新卡片 `例8.5（续）` 承接，并在 PDF p96（85）排完
%     （含法二、法三及蓝色模板表）；
%   · 例 8.6 题干在 PDF p97（86），其【解】(1) 跨 PDF p97 末～p98（87），
%     (2) 的解答在 PDF p99（印刷页 88，属下一分片 ch08_c）续完（该片以正文承接，未另开卡片），
%     故本片在 p98 末结束 fcard。
% ===========================================================

% -----------------------------------------------------------
% PDF p95（印刷页 84）—— 例 8.5 续（法一收尾、法二）
% -----------------------------------------------------------
% 承 PDF p94（印刷页 83）例 8.5【解】法一：对 $[P\,\vdots\,E]$ 作初等行变换求 $P^{-1}$
\begin{fcard}{例8.5（续）}
\fml{\to\left[\begin{array}{ccc|ccc}1&0&0&-\frac13&\frac23&-\frac13\\0&1&0&-\frac13&-\frac13&\frac23\\0&0&1&\frac13&\frac13&\frac13\end{array}\right],}
则 $P^{-1}=\frac13\begin{pmatrix}-1&2&-1\\-1&-1&2\\1&1&1\end{pmatrix}$，故
\fml{A^{n}=PA^{n}P^{-1}=\frac13\begin{pmatrix}-1&-1&1\\1&0&1\\0&1&1\end{pmatrix}\begin{pmatrix}1&&\\&1&\\&&3^{n}\end{pmatrix}\begin{pmatrix}-1&2&-1\\-1&-1&2\\1&1&1\end{pmatrix}}
\fml{=\frac13\begin{pmatrix}-1&-1&1\\1&0&1\\0&1&1\end{pmatrix}\begin{pmatrix}1&&\\&1&\\&&3^{n}\end{pmatrix}\begin{pmatrix}-1&2&-1\\-1&-1&2\\1&1&1\end{pmatrix}}
\fml{=\frac13\begin{pmatrix}2+3^{n}&-1+3^{n}&-1+3^{n}\\-1+3^{n}&2+3^{n}&-1+3^{n}\\-1+3^{n}&-1+3^{n}&2+3^{n}\end{pmatrix}.}

法二\quad 由法一得，$A\xi_3=\lambda_3\xi_3$，其中 $\lambda_3=3$，$\xi_3=[1,1,1]^{\mathrm T}$. 所以对应于 $\lambda_3=3$ 的全部特征向量为 $k_3\xi_3$（$k_3$ 为任意非零常数）.

设 $\lambda_1=\lambda_2=1$ 对应的特征向量为 $\xi=[x_1,x_2,x_3]^{\mathrm T}$，则
\fml{\xi_3^{\mathrm T}\xi=x_1+x_2+x_3=0,}
取 $\xi_1=[1,-1,0]^{\mathrm T}$，再取 $\xi_2$ 与 $\xi_1$ 正交，设 $\xi_2=[1,1,x]^{\mathrm T}$，代入上式得 $\xi_2=[1,1,-2]^{\mathrm T}$，所以对应于 $\lambda_1=\lambda_2=1$ 的全部特征向量为 $k_1\xi_1+k_2\xi_2$（$k_1$，$k_2$ 为不全为零的任意常数）.

将 $\xi_1$，$\xi_2$，$\xi_3$ 单位化，并取正交矩阵
\fml{Q=[\xi_1^{*},\xi_2^{*},\xi_3^{*}]=\begin{pmatrix}\frac{1}{\sqrt2}&\frac{1}{\sqrt6}&\frac{1}{\sqrt3}\\-\frac{1}{\sqrt2}&\frac{1}{\sqrt6}&\frac{1}{\sqrt3}\\0&-\frac{2}{\sqrt6}&\frac{1}{\sqrt3}\end{pmatrix},}
% 数学操作旁注（原稿蓝色手写）：$Q$ 下方箭头 → $A=Q\Lambda Q^{-1}$；$Q^{-1}$ 下方标注 $Q^{\mathrm T}$（即 $Q^{-1}=Q^{\mathrm T}$）
（旁注：$A=Q\Lambda Q^{-1}$，其中 $Q^{-1}=Q^{\mathrm T}$）
则
\fml{Q^{-1}AQ=Q^{\mathrm T}AQ=\begin{pmatrix}1&&\\&1&\\&&3\end{pmatrix}=\Lambda,}

% -----------------------------------------------------------
% PDF p96（印刷页 85）—— 例 8.5 续（法二收尾、法三）；【注】(1)(2)
% -----------------------------------------------------------
\fml{A^{n}=Q\Lambda^{n}Q^{\mathrm T}=\begin{pmatrix}\frac{1}{\sqrt2}&\frac{1}{\sqrt6}&\frac{1}{\sqrt3}\\-\frac{1}{\sqrt2}&\frac{1}{\sqrt6}&\frac{1}{\sqrt3}\\0&-\frac{2}{\sqrt6}&\frac{1}{\sqrt3}\end{pmatrix}\begin{pmatrix}1&&\\&1&\\&&3^{n}\end{pmatrix}\begin{pmatrix}\frac{1}{\sqrt2}&-\frac{1}{\sqrt2}&0\\\frac{1}{\sqrt6}&\frac{1}{\sqrt6}&-\frac{2}{\sqrt6}\\\frac{1}{\sqrt3}&\frac{1}{\sqrt3}&\frac{1}{\sqrt3}\end{pmatrix}=\frac13\begin{pmatrix}2+3^{n}&-1+3^{n}&-1+3^{n}\\-1+3^{n}&2+3^{n}&-1+3^{n}\\-1+3^{n}&-1+3^{n}&2+3^{n}\end{pmatrix}.}

法三\quad 由法一得，$A$ 的特征值为 $1,1,3$，$A$ 对应于 $\lambda_3=3$ 的特征向量为 $\xi_3=[1,1,1]^{\mathrm T}$，则 $A^{n}$ 的特征值为 $1,1,3^{n}$，$A^{n}-E$ 的特征值为 $0,0,3^{n}-1$，其对应于 $3^{n}-1$ 的特征向量仍为 $\xi_3=[1,1,1]^{\mathrm T}$，

单位化得 $\eta=\frac{1}{\sqrt3}[1,1,1]^{\mathrm T}=\frac{1}{\sqrt3}\xi_3$.

从而由实对称矩阵的相关结论（见注（2））得
\fml{A^{n}-E=(3^{n}-1)\eta\eta^{\mathrm T}=(3^{n}-1)\frac{1}{\sqrt3}\xi_3\times\frac{1}{\sqrt3}\xi_3^{\mathrm T}=\frac13(3^{n}-1)\begin{pmatrix}1&1&1\\1&1&1\\1&1&1\end{pmatrix},}
故 $A^{n}=\frac13(3^{n}-1)\begin{pmatrix}1&1&1\\1&1&1\\1&1&1\end{pmatrix}+E=\frac13\begin{pmatrix}2+3^{n}&-1+3^{n}&-1+3^{n}\\-1+3^{n}&2+3^{n}&-1+3^{n}\\-1+3^{n}&-1+3^{n}&2+3^{n}\end{pmatrix}$.

% 数学旁注（原稿蓝色手写模板表）
\fml{\begin{array}{c|c}A&f(A)\\\hline\lambda&f(\lambda)\\\hline\xi&\xi\end{array}}

% 数学旁注（原稿蓝色手写，接上表）：
$A$ 的特征值为 $1,1,3$；\\
$A^{n}$ 的特征值为 $1,1,3^{n}$；\\
$A^{n}-E$ 的特征值为 $0,0,3^{n}-1$；\\
$A^{n}-E=(3^{n}-1)\eta\eta^{\mathrm T}$；\\
$A^{n}=(3^{n}-1)\eta\eta^{\mathrm T}+E$
\end{fcard}

\begin{fnote}
【注】（1）因为 $A$ 是实对称矩阵，所以不同特征值对应的特征向量正交，且不仅存在可逆矩阵 $P$，使 $P^{-1}AP=\Lambda$，还存在正交矩阵 $Q$，使
\fml{Q^{-1}AQ=Q^{\mathrm T}AQ=\Lambda.}
法二中利用正交矩阵 $Q$，有 $Q^{-1}=Q^{\mathrm T}$，避免了用初等变换求逆矩阵，较简便.

（2）设 $n$ 阶实对称矩阵 $A$ 属于特征值 $\lambda_1,\lambda_2,\dots,\lambda_n$ 的单位正交特征向量为 $\xi_1,\xi_2,\dots,\xi_n$，则
\fml{A=\lambda_1\xi_1\xi_1^{\mathrm T}+\lambda_2\xi_2\xi_2^{\mathrm T}+\dots+\lambda_n\xi_n\xi_n^{\mathrm T}.}
证\quad 令 $Q=[\xi_1,\xi_2,\dots,\xi_n]$，有 $Q^{\mathrm T}AQ=\Lambda$，即
\fml{A=Q\Lambda Q^{\mathrm T}=[\xi_1,\xi_2,\dots,\xi_n]\begin{pmatrix}\lambda_1&&&\\&\lambda_2&&\\&&\ddots&\\&&&\lambda_n\end{pmatrix}\begin{pmatrix}\xi_1^{\mathrm T}\\\xi_2^{\mathrm T}\\\vdots\\\xi_n^{\mathrm T}\end{pmatrix}}
\fml{=\lambda_1\xi_1\xi_1^{\mathrm T}+\lambda_2\xi_2\xi_2^{\mathrm T}+\dots+\lambda_n\xi_n\xi_n^{\mathrm T}.}
\end{fnote}

% 数学旁注（原稿蓝色手写，【注】框右下方引出）：
% → 若 $\lambda_1=100$，$0\le\lambda_2,\lambda_3,\dots,\lambda_n<0.1$，你是否看出 $A$ 的主要成分是 $\lambda_1\xi_1\xi_1^{\mathrm T}$？
若 $\lambda_1=100$，$0\le\lambda_2,\lambda_3,\dots,\lambda_n<0.1$，你是否看出 $A$ 的主要成分是 $\lambda_1\xi_1\xi_1^{\mathrm T}$？

% -----------------------------------------------------------
% PDF p97（印刷页 86）
% -----------------------------------------------------------
% OPD 蓝色标注（原稿）：节标题后附「（$O_4$（盯住目标 4）$+D_1$（常规操作）$+D_{21}$（观察研究对象）$+P_4$（逆否思路））」——按 AGENTS.md §7.1 统一口径整段删除
% OPD 蓝色标注（原稿）：版心左侧竖排模块名「隐含条件体系块」（OPD 体系块标记，已删）
% 原稿右上角：蓝色装饰性二维码 1 个（非正文，未转写）

\fsec{四、$A$ 与 $B$ 相似}

\begin{fcard}{（1）}
若 $A$ 相似于 $B$，则\\
①$|A|=|B|$；\\
②$r(A)=r(B)$；\\
③$\mathrm{tr}(A)=\mathrm{tr}(B)$；\\
④$\lambda_A=\lambda_B$（或 $|\lambda E-A|=|\lambda E-B|$）；\\
⑤属于 $\lambda_A$ 的线性无关的特征向量的个数等于属于 $\lambda_B$ 的线性无关的特征向量的个数；\\
⑥$A$，$B$ 的各阶主子式之和分别相等.
% OPD 蓝色标注（原稿）：首句「则」右侧 →$D_{22}$（转换等价表述），当一个结论有"充要""充分""必要""否定"等条件，从而构成了一个好的命题点"时，便给考生指出了重点复习的方向
\end{fcard}

\begin{fcard}{（2）}
若 $A$ 相似于 $\Lambda$，$B$ 相似于 $\Lambda$，则 $A$ 相似于 $B$.

% 数学旁注（原稿蓝色手写证明）：
求可逆矩阵 $P$，使 $P^{-1}AP=B$；$\exists P_1$，$P_1^{-1}AP_1=\Lambda$；$\exists P_2$，$P_2^{-1}BP_2=\Lambda\Rightarrow P_1^{-1}AP_1=P_2^{-1}BP_2\Rightarrow$
\fml{P_2P_1^{-1}AP_1P_2^{-1}=B}
（原稿在 $P_2P_1^{-1}$ 下方标注 $P^{-1}$，在 $P_1P_2^{-1}$ 下方标注 $P$）
\end{fcard}

\begin{fcard}{（3）}
若 $A$ 相似于 $B$，$B$ 相似于 $\Lambda$，则 $A$ 相似于 $\Lambda$.

% 数学旁注（原稿蓝色手写证明）：
求可逆矩阵 $P$，使 $P^{-1}AP=\Lambda$；$\exists P_1$，$P_1^{-1}AP_1=B$，$\exists P_2$，$P_2^{-1}BP_2=\Lambda$，即 $B=P_2\Lambda P_2^{-1}\Rightarrow P_1^{-1}AP_1=P_2\Lambda P_2^{-1}\Rightarrow$
\fml{P_2^{-1}P_1^{-1}AP_1P_2=\Lambda}
（原稿在 $P_2^{-1}P_1^{-1}$ 下方标注 $P^{-1}$，在 $P_1P_2$ 下方标注 $P$）
\end{fcard}

\begin{fcard}{（4）}
$A$ 与 $B$ 相似手段的"三同一不同".

若 $P^{-1}AP=B$，则 $P^{-1}f(A)P=f(B)$，$P^{-1}A^{-1}P=B^{-1}$，$P^{-1}A^{*}P=B^{*}$，即 $f(A)$ 与 $f(B)$，$A^{-1}$ 与 $B^{-1}$，$A^{*}$ 与 $B^{*}$ 相似的手段相同，也即 $P^{-1}[af(A)+bA^{-1}+cA^{*}]P=af(B)+bB^{-1}+cB^{*}$. 但 $A^{\mathrm T}$ 与 $B^{\mathrm T}$ 相似的手段与上面不同.
\end{fcard}

\begin{fcard}{★★★例8.6}
若矩阵 $A$ 的伴随矩阵 $A^{*}=\begin{pmatrix}3&2&-2\\0&-1&0\\a&2&-3\end{pmatrix}$ 相似于矩阵 $B=\begin{pmatrix}-1&0&0\\-2&1&0\\0&0&-1\end{pmatrix}$，其中 $|A|>0$.

（1）求 $a$ 的值，并求可逆矩阵 $Q$，使得 $Q^{-1}A^{*}Q=B$；

（2）求 $A^{99}$.

【解】（1）因为矩阵 $A^{*}$ 与矩阵 $B$ 相似，所以它们的特征值相同. 又因为矩阵 $B$ 的特征多项式为
\fml{|\lambda E-B|=\begin{vmatrix}\lambda+1&0&0\\2&\lambda-1&0\\0&0&\lambda+1\end{vmatrix}=(\lambda+1)^2(\lambda-1),}
所以 $B$ 的特征值为 $\lambda_1=\lambda_2=-1$，$\lambda_3=1$.

由矩阵 $A^{*}$ 的特征多项式为

% -----------------------------------------------------------
% PDF p98（印刷页 87）—— 例 8.6 续
% -----------------------------------------------------------
\fml{|\lambda E-A^{*}|=\begin{vmatrix}\lambda-3&-2&2\\0&\lambda+1&0\\-a&-2&\lambda+3\end{vmatrix}=(\lambda+1)(\lambda^2-9+2a),}
则 $\lambda^2-9+2a=(\lambda+1)(\lambda-1)$，解得 $a=4$.

所以 $A^{*}=\begin{pmatrix}3&2&-2\\0&-1&0\\4&2&-3\end{pmatrix}$，设 $A^{*}$ 的特征值为 $\lambda^{*}$，且 $A^{*}$ 的特征值为 $\lambda_1^{*}=\lambda_2^{*}=-1$，$\lambda_3^{*}=1$，分别对应的特征向量为
\fml{\alpha_1=[1,-2,0]^{\mathrm T},\quad\alpha_2=[1,0,2]^{\mathrm T},\quad\alpha_3=[1,0,1]^{\mathrm T}.}

对 $B=\begin{pmatrix}-1&0&0\\-2&1&0\\0&0&-1\end{pmatrix}$，且 $B$ 的特征值为 $\lambda_1=\lambda_2=-1$，$\lambda_3=1$，分别对应的特征向量为
\fml{\beta_1=[1,1,0]^{\mathrm T},\quad\beta_2=[0,0,1]^{\mathrm T},\quad\beta_3=[0,1,0]^{\mathrm T}.}

则有
\fml{P_1^{-1}A^{*}P_1=\Lambda=\begin{pmatrix}-1&&\\&-1&\\&&1\end{pmatrix},\quad P_2^{-1}BP_2=\Lambda=\begin{pmatrix}-1&&\\&-1&\\&&1\end{pmatrix},}
其中 $P_1=[\alpha_1,\alpha_2,\alpha_3]=\begin{pmatrix}1&1&1\\-2&0&0\\0&2&1\end{pmatrix}$，$P_2=[\beta_1,\beta_2,\beta_3]=\begin{pmatrix}1&0&0\\1&0&1\\0&1&0\end{pmatrix}$，故有 $P_2P_1^{-1}A^{*}P_1P_2^{-1}=B$，即 $(P_1P_2^{-1})^{-1}A^{*}P_1P_2^{-1}=B$，故存在可逆矩阵
\fml{Q=P_1P_2^{-1}=\begin{pmatrix}1&1&1\\-2&0&0\\0&2&1\end{pmatrix}\begin{pmatrix}1&0&0\\1&0&1\\0&1&0\end{pmatrix}^{-1}=\begin{pmatrix}1&1&1\\-2&0&0\\0&2&1\end{pmatrix}\begin{pmatrix}1&0&0\\0&0&1\\-1&1&0\end{pmatrix}}
\fml{=\begin{pmatrix}0&1&1\\-2&0&0\\-1&1&2\end{pmatrix},}
使 $Q^{-1}A^{*}Q=B$.

（2）由（1）知，$|A^{*}|\ne0$，故 $A^{*}$ 可逆，$\lambda^{*}$ 为矩阵 $A^{*}$ 的特征值，对应的特征向量为 $\alpha$，则有 $A^{*}\alpha=\lambda^{*}\alpha$，

等式两端的左边乘矩阵 $A$ 得 $AA^{*}\alpha=\lambda^{*}A\alpha$，即 $A\alpha=\frac{|A|}{\lambda^{*}}\alpha$，故 $\frac{|A|}{\lambda^{*}}$ 是矩阵 $A$ 的特征值，$\alpha$ 是矩阵 $A$ 的对应于特征值 $\frac{|A|}{\lambda^{*}}$ 的特征向量.
\fml{|A^{*}|=\begin{vmatrix}3&2&-2\\0&-1&0\\4&2&-3\end{vmatrix}=1=|A|^2,}
% 例 8.6 的（2）问解答在 PDF p99（印刷页 88，属下一分片）续完，此处结束 fcard
\end{fcard}
